\begin{table}[H]
\centering
\caption{Nama Daerah Perbandingan Cluster Biplot vs Kategori IPM (UNDP) (2024)} \label{tab:crossname2024}
\begin{adjustbox}{max width=\textwidth}
\begin{tabular}{|p{4cm}|p{10cm}|}
\hline
\textbf{Kategori} & \textbf{Daftar Wilayah} \\
\hline
Kl.1 \& Sangat Tinggi & - \\
\hline
Kl.1 \& Tinggi & Ciamis \\
\hline
Kl.1 \& Sedang & Cianjur, Tasikmalaya, Majalengka, Pangandaran, Pandeglang, Lebak \\
\hline
Kl.1 \& Rendah & - \\
\hline
Kl.2 \& Sangat Tinggi & - \\
\hline
Kl.2 \& Tinggi & Kep. Seribu, Bogor, Bandung, Kuningan, Cirebon, Sumedang, Subang, Purwakarta, Karawang, Bandung Barat, Kota Cirebon, Kota Tasikmalaya, Kota Banjar, Tangerang, Kota Cilegon, Kota Serang \\
\hline
Kl.2 \& Sedang & Sukabumi, Garut, Indramayu, Serang \\
\hline
Kl.2 \& Rendah & - \\
\hline
Kl.3 \& Sangat Tinggi & Kota Jakarta Selatan, Kota Jakarta Timur, Kota Jakarta Pusat, Kota Jakarta Barat, Kota Jakarta Utara, Kota Bandung, Kota Bekasi, Kota Depok, Kota Cimahi, Kota Tangerang, Kota Tangerang Selatan \\
\hline
Kl.3 \& Tinggi & Bekasi, Kota Bogor, Kota Sukabumi \\
\hline
Kl.3 \& Sedang & - \\
\hline
Kl.3 \& Rendah & - \\
\hline
\end{tabular}
\end{adjustbox}
\end{table}